\conclusion

В результате выполнения выпускной квалификационной работы бакалавра была разработана информационная система управления молочной фермой с интеграцией роботизированного доильного оборудования Lely и предиктивной аналитикой на основе моделей машинного обучения. Все поставленные задачи были выполнены.

Проведён анализ предметной области и существующих решений (Lely Horizon, DeLaval DelPro, SCR Heatime, 1С: Управление фермой), показавший, что ни одно из них не обеспечивает одновременно полный цикл управления стадом, интеграцию с оборудованием сторонних производителей, встроенные модели машинного обучения и открытую архитектуру. На основании анализа сформулировано техническое задание и обоснован выбор стека технологий (Rust/Axum, SvelteKit, Python/FastAPI, PostgreSQL/TimescaleDB).

Спроектирована и реализована трёхзвенная архитектура. Серверная часть на Rust/Axum включает 25 модулей обработчиков, 33 модуля сервисов, 24 модуля моделей данных и 6 модулей middleware. Аутентификация на основе JWT с HttpOnly cookie обеспечивает защиту от CSRF и XSS. Middleware-стек включает rate limiting (Redis/in-memory), gzip-сжатие, CSRF-защиту, метрики Prometheus и распределённую трассировку. База данных PostgreSQL содержит более 35 таблиц с расширением TimescaleDB для эффективной работы с временными рядами.

Реализована интеграция с облачным API Lely Horizon, обеспечивающая автоматическую двустороннюю синхронизацию 23 типов данных (животные, надои, воспроизводство, кормление, фитнес, оборудование) с шифрованием учётных данных AES-256-GCM и инкрементальной синхронизацией по временным диапазонам.

Клиентское веб-приложение на SvelteKit~5 содержит более 24 страниц с серверным рендерингом и адаптивным интерфейсом на Tailwind CSS. Разработано более 16 переиспользуемых UI-компонентов, включая DataTable с серверной пагинацией и адаптивным отображением, компоненты визуализации на Chart.js (надои, SCC, лактационные кривые, прогнозы), систему уведомлений и валидацию форм на основе реактивных Runes.

Модуль машинного обучения на Python/FastAPI включает 11 моделей: прогноз надоев на корову (XGBoost с квантильной регрессией q10/q50/q90, обеспечивающей доверительные интервалы) и на стадо (Prophet с декомпозицией тренда и сезонности), обнаружение мастита и кетоза, выявление охоты, оценка риска выбраковки, оценка BCS, кластеризация стада (KMeans с автоматическим выбором $k$ по силуэту), аномалии оборудования (Isolation Forest) и рекомендации по кормлению. MLOps-инфраструктура включает: автоматическое переобучение (APScheduler), оптимизацию гиперпараметров (Optuna, TPE-сэмплер, 8-мерное пространство), мониторинг дрейфа данных (prediction drift и feature drift), версионирование моделей и отслеживание экспериментов (MLflow).

В серверной части реализованы алгоритмы предиктивной аналитики на Rust: метод Хольта--Винтерса для прогнозирования временных рядов надоя с обнаружением выбросов (IQR) и структурных сдвигов (t-критерий); кривая лактации Вуда для моделирования зависимости надоя от DIM; комплексный индекс здоровья на основе Z-оценок четырёх показателей (надой, жвачка, активность, SCC) с динамической визуализацией тренда за 90 дней; оптимизатор запуска; система сравнения 8 моделей прогнозирования (SES, SMA, WMA, линейная регрессия, метод Хольта, Хольт--Винтерс, ряд Фурье, ARIMA(0,1,1)) с автоматическим выбором лучшей по MAPE. Реализован ансамблевый прогноз, объединяющий ML-модель (XGBoost) и лучшую Rust-модель методом взвешенного среднего с весами, обратно пропорциональными MAPE, что повышает робастность прогнозирования. Для обеспечения интерпретируемости прогнозов внедрён метод SHAP, визуализирующий вклад каждого признака в виде горизонтальной столбчатой диаграммы в интерфейсе.

Основные количественные результаты разработки включают: серверную часть из 25 модулей обработчиков, 33 модулей сервисов, 24 модулей моделей данных и 6 модулей middleware; базу данных из более 35 таблиц, более 20 файлов миграций с расширениями TimescaleDB и pg\_trgm; клиентское приложение из более 24 страниц, 23 модулей API-клиента и более 16 UI-компонентов; модуль машинного обучения из 11 моделей с автоматическим обучением и переобучением, мониторингом дрейфа и оптимизацией гиперпараметров через Optuna; серверную аналитику из 8 моделей прогнозирования с автоматическим сравнением по MAPE/RMSE, индексом здоровья и кривой Вуда; интеграцию с Lely, охватывающую 23 типа синхронизируемых данных с шифрованием AES-256-GCM и инкрементальной синхронизацией; мониторинг через стек Grafana LGTM для сбора метрик, логов и трассировок; тестирование из 30+ тестовых файлов с нагрузочным тестированием через k6; и конвейер CI/CD с параллельным тестированием трёх компонентов и автоматическим деплоем при создании тегов.

Научная новизна работы заключается в архитектурном решении, позволяющем серверным алгоритмам прогнозирования на Rust работать автономно, обеспечивая предиктивную аналитику даже при недоступности ML-сервиса, а также в реализации системы автоматического сравнения 8 моделей прогнозирования с выбором оптимальной для каждого животного индивидуально и ансамблевого объединения прогнозов ML и Rust-моделей с весами, обратно пропорциональными MAPE.

Практическая значимость подтверждается полнотой покрытия бизнес-процессов молочной фермы: от учёта животных и надоев до предиктивной аналитики и интеграции с роботизированным оборудованием. Система развёрнута в контейнерах Docker с обратным прокси Nginx, SSL-терминацией и мониторингом, что обеспечивает готовность к промышленной эксплуатации.

Дальнейшее развитие системы может включать: расширение набора моделей машинного обучения (прогнозирование заболеваемости конечностей, оптимизация параметров доильных роботов); интеграцию с дополнительными производителями оборудования (DeLaval, GEA); разработку мобильного приложения для фермеров с push-уведомлениями; добавление модулей финансового анализа и планирования с учётом стоимости кормов, ветеринарных услуг и рыночных цен на молоко; внедрение online learning для адаптации моделей к новым данным в реальном времени.
